% ORIG03-SEGMENT-CODE: OR03-S004
% ORIG03-SEGMENT-ID: d90.orig.v1.tr03.seg0004
% ORIG03-MODULE-ID: ps02.module.0001
% SPDX-License-Identifier: CC-BY-SA-4.0
% Karya asli berbahasa Indonesia; semua butir disusun independen.

\section{Set soal II: metode proksimal}
\label{orig03:section:set-soal-metode-proksimal}

Lima masalah berikut menghubungkan operator proksimal, pemisahan
maju--mundur, lemma penurunan, resolven monoton, dan laju titik proksimal.

% ORIG03-STABLE-ID: ps02.problem.0001
\begin{exercise}[Ambang lunak sebagai operator proksimal]
\label{orig03:ps02:problem:0001}
Untuk $\lambda>0$, definisikan
$S_\lambda(v)=\operatorname{sign}(v)\max\{|v|-\lambda,0\}$.
\begin{enumerate}
% ORIG03-STABLE-ID: ps02.prompt.0001
\item\label{orig03:ps02:prompt:0001}
Buktikan bahwa $\operatorname{prox}_{\lambda|\cdot|}(v)=S_\lambda(v)$ untuk
setiap $v\in\mathbb{R}$.
% ORIG03-STABLE-ID: ps02.prompt.0002
\item\label{orig03:ps02:prompt:0002}
Dengan penerapan per koordinat, hitung
$\operatorname{prox}_{(3/4)\lVert\cdot\rVert_1}(3,-1/2,-2)$.
\end{enumerate}
\end{exercise}

% ORIG03-STABLE-ID: ps02.hint1.0001
\par\noindent\textbf{Petunjuk tahap 1 (ps02.prompt.0001).}
Tuliskan syarat optimalitas untuk
$\min_x\{\lambda|x|+\tfrac12(x-v)^2\}$.
\label{orig03:ps02:hint1:0001}
% ORIG03-STABLE-ID: ps02.hint2.0001
\par\noindent\textbf{Petunjuk tahap 2 (ps02.prompt.0001).}
Pisahkan kasus $x>0$, $x<0$, dan $x=0$ dalam
$0\in x-v+\lambda\partial|x|$.
\label{orig03:ps02:hint2:0001}
% ORIG03-STABLE-ID: ps02.answer.0001
\par\noindent\textbf{Jawaban ringkas (ps02.prompt.0001).}
\[
 \operatorname{prox}_{\lambda|\cdot|}(v)=
 \begin{cases}
 v-\lambda,&v>\lambda,\\
 0,&|v|\leq\lambda,\\
 v+\lambda,&v<-\lambda,
 \end{cases}
 =S_\lambda(v).
\]
\label{orig03:ps02:answer:0001}
% ORIG03-STABLE-ID: ps02.solution.0001
\par\noindent\textbf{Solusi lengkap (ps02.prompt.0001).}
Objektif bersifat konveks kuat, sehingga syarat
\[
 0\in x-v+\lambda\partial|x|
\]
menentukan peminimum unik. Jika $x>0$, syaratnya $x=v-\lambda$, konsisten
tepat ketika $v>\lambda$. Jika $x<0$, diperoleh $x=v+\lambda$, konsisten
tepat ketika $v<-\lambda$. Jika $x=0$, syarat menjadi
$v\in\lambda[-1,1]$, ekuivalen dengan $|v|\leq\lambda$. Ketiga kasus itu
persis rumus $S_\lambda(v)$.
\emph{Rubrik bukti: \ref{orig03:rubric:proof:0005}.}
\label{orig03:ps02:solution:0001}

% ORIG03-STABLE-ID: ps02.hint1.0002
\par\noindent\textbf{Petunjuk tahap 1 (ps02.prompt.0002).}
Norma satu dan kuadrat jarak keduanya terpisah menurut koordinat.
\label{orig03:ps02:hint1:0002}
% ORIG03-STABLE-ID: ps02.hint2.0002
\par\noindent\textbf{Petunjuk tahap 2 (ps02.prompt.0002).}
Terapkan $S_{3/4}$ pada ketiga koordinat.
\label{orig03:ps02:hint2:0002}
% ORIG03-STABLE-ID: ps02.answer.0002
\par\noindent\textbf{Jawaban ringkas (ps02.prompt.0002).}
$\operatorname{prox}_{(3/4)\lVert\cdot\rVert_1}(3,-1/2,-2)
=(9/4,0,-5/4)$.
\label{orig03:ps02:answer:0002}
% ORIG03-STABLE-ID: ps02.solution.0002
\par\noindent\textbf{Solusi lengkap (ps02.prompt.0002).}
Pemisahan koordinat memberi
\[
 S_{3/4}(3)=3-\frac34=\frac94,
 \qquad S_{3/4}(-1/2)=0,
 \qquad S_{3/4}(-2)=-2+\frac34=-\frac54.
\]
Menggabungkan ketiga peminimum skalar menghasilkan vektor yang dinyatakan.
\label{orig03:ps02:solution:0002}

% ORIG03-STABLE-ID: ps02.problem.0002
\begin{exercise}[Iterasi maju--mundur skalar]
\label{orig03:ps02:problem:0002}
Pertimbangkan
$F(x)=\tfrac12(x-3)^2+|x|$ dan langkah $\gamma=1/2$.
\begin{enumerate}
% ORIG03-STABLE-ID: ps02.prompt.0003
\item\label{orig03:ps02:prompt:0003}
Turunkan pemetaan iterasi maju--mundur $x^+=T(x)$.
% ORIG03-STABLE-ID: ps02.prompt.0004
\item\label{orig03:ps02:prompt:0004}
Dari $x_0=0$, hitung $x_1$ dan $x_2$, lalu tentukan peminimum eksak $F$.
\end{enumerate}
\end{exercise}

% ORIG03-STABLE-ID: ps02.hint1.0003
\par\noindent\textbf{Petunjuk tahap 1 (ps02.prompt.0003).}
Pisahkan bagian mulus $g(x)=\tfrac12(x-3)^2$ dan $h(x)=|x|$.
\label{orig03:ps02:hint1:0003}
% ORIG03-STABLE-ID: ps02.hint2.0003
\par\noindent\textbf{Petunjuk tahap 2 (ps02.prompt.0003).}
Gunakan $T(x)=\operatorname{prox}_{\gamma h}(x-\gamma g'(x))$.
\label{orig03:ps02:hint2:0003}
% ORIG03-STABLE-ID: ps02.answer.0003
\par\noindent\textbf{Jawaban ringkas (ps02.prompt.0003).}
$T(x)=S_{1/2}((x+3)/2)$.
\label{orig03:ps02:answer:0003}
% ORIG03-STABLE-ID: ps02.solution.0003
\par\noindent\textbf{Solusi lengkap (ps02.prompt.0003).}
Karena $g'(x)=x-3$ dan $\gamma=1/2$,
\[
 x-\gamma g'(x)=x-\frac12(x-3)=\frac{x+3}{2}.
\]
Langkah mundur adalah operator proksimal $\operatorname{prox}_{(1/2)|\cdot|}
=S_{1/2}$, sehingga $T(x)=S_{1/2}((x+3)/2)$.
\label{orig03:ps02:solution:0003}

% ORIG03-STABLE-ID: ps02.hint1.0004
\par\noindent\textbf{Petunjuk tahap 1 (ps02.prompt.0004).}
Substitusikan berturut-turut $x_0$ dan $x_1$ ke $T$.
\label{orig03:ps02:hint1:0004}
% ORIG03-STABLE-ID: ps02.hint2.0004
\par\noindent\textbf{Petunjuk tahap 2 (ps02.prompt.0004).}
Untuk peminimum, selesaikan $0\in x-3+\partial|x|$.
\label{orig03:ps02:hint2:0004}
% ORIG03-STABLE-ID: ps02.answer.0004
\par\noindent\textbf{Jawaban ringkas (ps02.prompt.0004).}
$x_1=1$, $x_2=3/2$, dan peminimum uniknya $x^\star=2$.
\label{orig03:ps02:answer:0004}
% ORIG03-STABLE-ID: ps02.solution.0004
\par\noindent\textbf{Solusi lengkap (ps02.prompt.0004).}
Kita memperoleh
\[
 x_1=S_{1/2}(3/2)=1,
 \qquad
 x_2=S_{1/2}(2)=\frac32.
\]
Syarat optimalitas adalah $0\in x-3+\partial|x|$. Pada cabang $x>0$,
syarat itu menjadi $0=x-2$, sehingga $x=2$, yang konsisten dengan cabang.
Cabang $x<0$ tidak menghasilkan titik yang konsisten, dan di nol interval
$-3+[-1,1]$ tidak memuat nol. Konveksitas kuat suku kuadratik menjamin
keunikan, jadi $x^\star=2$.
\label{orig03:ps02:solution:0004}

% ORIG03-STABLE-ID: ps02.problem.0003
\begin{exercise}[Lemma penurunan]
\label{orig03:ps02:problem:0003}
Misalkan $f:\mathbb{R}^n\to\mathbb{R}$ terdiferensialkan dan gradiennya
$L$-Lipschitz.
\begin{enumerate}
% ORIG03-STABLE-ID: ps02.prompt.0005
\item\label{orig03:ps02:prompt:0005}
Buktikan
\[
 f(y)\leq f(x)+\nabla f(x)^\top(y-x)+\frac L2\lVert y-x\rVert_2^2.
\]
% ORIG03-STABLE-ID: ps02.prompt.0006
\item\label{orig03:ps02:prompt:0006}
Untuk $f(x)=\tfrac L2x^2$ pada $\mathbb{R}$, terapkan langkah gradien
$x^+=x-L^{-1}f'(x)$ dan periksa ketajaman batas penurunan yang dihasilkan.
\end{enumerate}
\end{exercise}

% ORIG03-STABLE-ID: ps02.hint1.0005
\par\noindent\textbf{Petunjuk tahap 1 (ps02.prompt.0005).}
Integralkan gradien sepanjang ruas dari $x$ ke $y$.
\label{orig03:ps02:hint1:0005}
% ORIG03-STABLE-ID: ps02.hint2.0005
\par\noindent\textbf{Petunjuk tahap 2 (ps02.prompt.0005).}
Dengan $d=y-x$, batasi
$[\nabla f(x+td)-\nabla f(x)]^\top d$ oleh $Lt\lVert d\rVert_2^2$.
\label{orig03:ps02:hint2:0005}
% ORIG03-STABLE-ID: ps02.answer.0005
\par\noindent\textbf{Jawaban ringkas (ps02.prompt.0005).}
Batas mengikuti teorema dasar kalkulus pada $t\mapsto f(x+t(y-x))$ dan
integral $\int_0^1Lt\,dt=L/2$.
\label{orig03:ps02:answer:0005}
% ORIG03-STABLE-ID: ps02.solution.0005
\par\noindent\textbf{Solusi lengkap (ps02.prompt.0005).}
Tetapkan $d=y-x$. Teorema dasar kalkulus memberi
\[
 f(y)-f(x)=\int_0^1\nabla f(x+td)^\top d\,dt
 =\nabla f(x)^\top d+
 \int_0^1[\nabla f(x+td)-\nabla f(x)]^\top d\,dt.
\]
Cauchy--Schwarz dan sifat Lipschitz menghasilkan
\[
 [\nabla f(x+td)-\nabla f(x)]^\top d
 \leq \lVert\nabla f(x+td)-\nabla f(x)\rVert_2\lVert d\rVert_2
 \leq Lt\lVert d\rVert_2^2.
\]
Mengintegralkan batas terakhir dari nol sampai satu membuktikan klaim.
\emph{Rubrik bukti: \ref{orig03:rubric:proof:0002}.}
\label{orig03:ps02:solution:0005}

% ORIG03-STABLE-ID: ps02.hint1.0006
\par\noindent\textbf{Petunjuk tahap 1 (ps02.prompt.0006).}
Hitung $f'(x)=Lx$.
\label{orig03:ps02:hint1:0006}
% ORIG03-STABLE-ID: ps02.hint2.0006
\par\noindent\textbf{Petunjuk tahap 2 (ps02.prompt.0006).}
Bandingkan hasil dengan
$f(x^+)\leq f(x)-(2L)^{-1}|f'(x)|^2$.
\label{orig03:ps02:hint2:0006}
% ORIG03-STABLE-ID: ps02.answer.0006
\par\noindent\textbf{Jawaban ringkas (ps02.prompt.0006).}
$x^+=0$ dan
$f(x^+)=f(x)-(2L)^{-1}|f'(x)|^2=0$; batasnya mencapai kesamaan.
\label{orig03:ps02:answer:0006}
% ORIG03-STABLE-ID: ps02.solution.0006
\par\noindent\textbf{Solusi lengkap (ps02.prompt.0006).}
Karena $f'(x)=Lx$, langkah berukuran $1/L$ memberi $x^+=x-x=0$.
Lemma penurunan dengan $y=x-L^{-1}f'(x)$ memberi
\[
 f(x^+)\leq f(x)-\frac1{2L}|f'(x)|^2.
\]
Untuk kuadratik ini, ruas kanan sama dengan
$\tfrac L2x^2-\tfrac1{2L}L^2x^2=0=f(0)=f(x^+)$, sehingga batas tersebut
tajam untuk setiap $x$.
\label{orig03:ps02:solution:0006}

% ORIG03-STABLE-ID: ps02.problem.0004
\begin{exercise}[Pemisahan maju--mundur sebagai kontraksi]
\label{orig03:ps02:problem:0004}
Pada $\mathbb{R}$, ambil $A=\partial|\cdot|$, $B(x)=x-2$, dan $\gamma=1/2$.
\begin{enumerate}
% ORIG03-STABLE-ID: ps02.prompt.0007
\item\label{orig03:ps02:prompt:0007}
Hitung $T=J_{\gamma A}(I-\gamma B)$ secara eksplisit.
% ORIG03-STABLE-ID: ps02.prompt.0008
\item\label{orig03:ps02:prompt:0008}
Tentukan titik tetap $T$ dan hubungkan dengan peminimum
$|x|+\tfrac12(x-2)^2$.
% ORIG03-STABLE-ID: ps02.prompt.0009
\item\label{orig03:ps02:prompt:0009}
Buktikan bahwa $T$ adalah kontraksi dengan faktor paling besar $1/2$ dan
simpulkan keunikan titik tetapnya.
\end{enumerate}
\end{exercise}

% ORIG03-STABLE-ID: ps02.hint1.0007
\par\noindent\textbf{Petunjuk tahap 1 (ps02.prompt.0007).}
$J_{\gamma\partial|\cdot|}=\operatorname{prox}_{\gamma|\cdot|}$.
\label{orig03:ps02:hint1:0007}
% ORIG03-STABLE-ID: ps02.hint2.0007
\par\noindent\textbf{Petunjuk tahap 2 (ps02.prompt.0007).}
Hitung $x-\tfrac12(x-2)=x/2+1$.
\label{orig03:ps02:hint2:0007}
% ORIG03-STABLE-ID: ps02.answer.0007
\par\noindent\textbf{Jawaban ringkas (ps02.prompt.0007).}
$T(x)=S_{1/2}(x/2+1)$.
\label{orig03:ps02:answer:0007}
% ORIG03-STABLE-ID: ps02.solution.0007
\par\noindent\textbf{Solusi lengkap (ps02.prompt.0007).}
Definisi resolven dan identitas resolven subdiferensial memberi
\[
 T(x)=\operatorname{prox}_{(1/2)|\cdot|}
 \left(x-\frac12(x-2)\right)
 =S_{1/2}\left(\frac x2+1\right).
\]
\label{orig03:ps02:solution:0007}

% ORIG03-STABLE-ID: ps02.hint1.0008
\par\noindent\textbf{Petunjuk tahap 1 (ps02.prompt.0008).}
Coba $x=1$ dalam rumus eksplisit $T$.
\label{orig03:ps02:hint1:0008}
% ORIG03-STABLE-ID: ps02.hint2.0008
\par\noindent\textbf{Petunjuk tahap 2 (ps02.prompt.0008).}
Bandingkan persamaan titik tetap dengan
$0\in x-2+\partial|x|$.
\label{orig03:ps02:hint2:0008}
% ORIG03-STABLE-ID: ps02.answer.0008
\par\noindent\textbf{Jawaban ringkas (ps02.prompt.0008).}
Titik tetapnya $x^\star=1$, yang juga merupakan peminimum unik objektif.
\label{orig03:ps02:answer:0008}
% ORIG03-STABLE-ID: ps02.solution.0008
\par\noindent\textbf{Solusi lengkap (ps02.prompt.0008).}
Rumus eksplisit memberi $T(1)=S_{1/2}(3/2)=1$. Selain itu,
$x=T(x)$ ekuivalen dengan
$x=J_{\gamma A}(x-\gamma Bx)$, lalu dengan
$0\in Ax+Bx=\partial|x|+x-2$. Ini persis syarat optimalitas untuk
$|x|+\tfrac12(x-2)^2$. Pada $x=1>0$, syaratnya $0=1+1-2$.
Suku kuadratik konveks kuat, jadi peminimum tersebut unik.
\label{orig03:ps02:solution:0008}

% ORIG03-STABLE-ID: ps02.hint1.0009
\par\noindent\textbf{Petunjuk tahap 1 (ps02.prompt.0009).}
Operator proksimal bersifat non-ekspansif.
\label{orig03:ps02:hint1:0009}
% ORIG03-STABLE-ID: ps02.hint2.0009
\par\noindent\textbf{Petunjuk tahap 2 (ps02.prompt.0009).}
Bandingkan argumen ambang lunak untuk $x$ dan $y$.
\label{orig03:ps02:hint2:0009}
% ORIG03-STABLE-ID: ps02.answer.0009
\par\noindent\textbf{Jawaban ringkas (ps02.prompt.0009).}
$|T(x)-T(y)|\leq\tfrac12|x-y|$; karena itu titik tetapnya tunggal.
\label{orig03:ps02:answer:0009}
% ORIG03-STABLE-ID: ps02.solution.0009
\par\noindent\textbf{Solusi lengkap (ps02.prompt.0009).}
Non-ekspansivitas $S_{1/2}$ menghasilkan
\[
 |T(x)-T(y)|
 \leq\left|\left(\frac x2+1\right)-
              \left(\frac y2+1\right)\right|
 =\frac12|x-y|.
\]
Jadi $T$ adalah kontraksi. Jika $x$ dan $y$ keduanya titik tetap, maka
$|x-y|=|T(x)-T(y)|\leq|x-y|/2$, yang memaksa $x=y$.
\emph{Rubrik bukti: \ref{orig03:rubric:proof:0005}.}
\label{orig03:ps02:solution:0009}

% ORIG03-STABLE-ID: ps02.problem.0005
\begin{exercise}[Laju titik proksimal untuk operator linear]
\label{orig03:ps02:problem:0005}
Ambil $A(x)=mx$ dengan $m>0$ dan parameter $\gamma>0$.
\begin{enumerate}
% ORIG03-STABLE-ID: ps02.prompt.0010
\item\label{orig03:ps02:prompt:0010}
Hitung resolven $J_{\gamma A}=(I+\gamma A)^{-1}$.
% ORIG03-STABLE-ID: ps02.prompt.0011
\item\label{orig03:ps02:prompt:0011}
Untuk iterasi titik proksimal $x_{k+1}=J_{\gamma A}(x_k)$, turunkan rumus
eksak $x_k$ dan lajunya menuju nol.
\end{enumerate}
\end{exercise}

% ORIG03-STABLE-ID: ps02.hint1.0010
\par\noindent\textbf{Petunjuk tahap 1 (ps02.prompt.0010).}
Tetapkan $z=J_{\gamma A}(x)$.
\label{orig03:ps02:hint1:0010}
% ORIG03-STABLE-ID: ps02.hint2.0010
\par\noindent\textbf{Petunjuk tahap 2 (ps02.prompt.0010).}
Selesaikan $x=z+\gamma mz$ terhadap $z$.
\label{orig03:ps02:hint2:0010}
% ORIG03-STABLE-ID: ps02.answer.0010
\par\noindent\textbf{Jawaban ringkas (ps02.prompt.0010).}
$J_{\gamma A}(x)=x/(1+\gamma m)$.
\label{orig03:ps02:answer:0010}
% ORIG03-STABLE-ID: ps02.solution.0010
\par\noindent\textbf{Solusi lengkap (ps02.prompt.0010).}
Definisi resolven menyatakan
$z=J_{\gamma A}(x)$ jika dan hanya jika
$x\in z+\gamma A(z)$. Karena $A(z)=mz$, ini menjadi
$x=(1+\gamma m)z$. Faktor $1+\gamma m$ positif, sehingga
$z=x/(1+\gamma m)$.
\emph{Rubrik bukti: \ref{orig03:rubric:proof:0005}.}
\label{orig03:ps02:solution:0010}

% ORIG03-STABLE-ID: ps02.hint1.0011
\par\noindent\textbf{Petunjuk tahap 1 (ps02.prompt.0011).}
Substitusikan rumus resolven ke relasi iterasi.
\label{orig03:ps02:hint1:0011}
% ORIG03-STABLE-ID: ps02.hint2.0011
\par\noindent\textbf{Petunjuk tahap 2 (ps02.prompt.0011).}
Ulangi perkalian oleh konstanta $(1+\gamma m)^{-1}$.
\label{orig03:ps02:hint2:0011}
% ORIG03-STABLE-ID: ps02.answer.0011
\par\noindent\textbf{Jawaban ringkas (ps02.prompt.0011).}
$x_k=(1+\gamma m)^{-k}x_0$ dan
$|x_k|=(1+\gamma m)^{-k}|x_0|\to0$ secara linear.
\label{orig03:ps02:answer:0011}
% ORIG03-STABLE-ID: ps02.solution.0011
\par\noindent\textbf{Solusi lengkap (ps02.prompt.0011).}
Rumus resolven memberi
$x_{k+1}=(1+\gamma m)^{-1}x_k$. Induksi terhadap $k$ menghasilkan
$x_k=(1+\gamma m)^{-k}x_0$. Karena $m,\gamma>0$, faktor
$q=(1+\gamma m)^{-1}$ berada dalam $(0,1)$. Maka
$|x_k|=q^k|x_0|$ dan konvergensinya menuju nol bersifat linear dengan faktor
eksak $q$.
\label{orig03:ps02:solution:0011}
